\chapter{Introduction}

In graph theory, the graph isomorphism problem (GI) is a well-known challenge: determining whether two given graphs are isomorphic. Fully retracing its history is nearly impossible given the sheer number of works devoted to the topic, but we will review its most significant milestones. GI first appeared as a computational problem in the chemical documentation literature around the 1950s~\cite{ray1957finding} (as cited in~\cite{grohe2020graph}), where it took the form of a matching problem between a ``molecular graph'' and graphs of the same type stored in a database. With this approach, we find ourselves at the origins of cheminformatics: the application of computer science to chemistry. Early instances of this matching problem gave rise to some of the first dedicated systems. Unger proposed GIT (\textbf{G}raph \textbf{I}somorphism \textbf{T}ester)~\cite{unger1964git}, one of the earliest heuristics for testing the isomorphism of pairs of directed line graphs, while Sussenguth introduced a graph-theoretic algorithm specifically designed for matching chemical structures~\cite{sussenguth1965graph}. These early, domain-specific efforts laid the foundations for the more general algorithmic developments that followed.

Throughout the 20th century, considerable research efforts were devoted to solving GI, with varying degrees of efficiency. On the one hand, one of the simplest approaches is based on backtracking, a widely known and relatively easy-to-implement technique. However, it suffers from poor scalability and provides no guarantee of efficiency in the worst case~\cite{schmidt1976fast}. On the other hand, more sophisticated methods were also proposed. One notable approach involved the construction of a so-called representative graph and a reordered graph to test for isomorphism~\cite{corneil1970efficient}. Although this general approach was later shown to be incorrect (as cited in~\cite{foggia2001performance}), it played an important role in stimulating further investigations into the graph isomorphism problem.

Finally, a more promising direction emerged: reducing graphs to their canonical form. The idea is straightforward: every graph can be transformed into a \textit{unique} canonical representation~\cite{arlazarov1974algorithm}, meaning that two isomorphic graphs will share the same canonical form. This shift in perspective transformed GI from a problem of structural equivalence into one of equality: rather than directly checking whether two graphs are isomorphic, one simply compares their canonical forms. This change of viewpoint led to more efficient solutions.

In 1981, Brendan D. McKay introduced a breakthrough with his software \Naut{} (``\textbf{n}o \textbf{aut}omorphisms, \textbf{y}es?'')~\cite{mckay1981practical}, which provided a much more efficient approach to tackling GI.\@ Initially, the program had been started in 1976, written in Fortran and Assembly, and was originally called GLABC.\@ However, in 1981, it was finalised in C and received its new name: \Naut{}. Since then, the program has undergone continuous optimisation to maximise its speed and efficiency. The key idea behind \Naut{} is to compute a graph's canonical form without relying on matrix representations. Instead, it uses a search tree combined with a partition refinement system to generate the canonical form.

Nowadays, this canonical form capability enables a powerful application: the exhaustive generation of non-isomorphic graphs. Given $n$ vertices (and optionally additional constraints), the goal is to enumerate all possible graphs that are pairwise non-isomorphic. Isomorph-free exhaustive generation is a problem that researchers have been studying for decades. For instance, the generation of cubic graphs has been investigated since the 1960s, as these graphs are useful in various fields: in chemistry, they can model carbon networks; in mathematics and computer science, they often serve as minimal counterexamples in graph theory~\cite{brinkmann2013history}. McKay, for his part, successfully connected his tool \Naut{} to the generation of non-isomorphic graphs~\cite{mckay1998isomorph}.

\Naut{} is much more than just a tool for canonical labelling and the generation of non-isomorphic graphs. It has evolved into a comprehensive program supporting a wide range of graph families and offering additional features, such as graph statistics and other advanced analysis tools~\cite{mckay2024nug}. Furthermore, libraries based on \Naut{} have been developed for several programming languages: Pynauty in Python~\cite{pynauty}, NautyGraphs in Julia~\cite{nautygraphs_jl}, and even pl-nauty in Prolog~\cite{frank2016logic}.

In practice, graph isomorphism tools are highly useful and widely used in scientific research. In general, these tools are employed to detect symmetries in combinatorial objects, which allows them to be applied across a variety of domains. In optimisation, for instance, GI tools are used in SAT solving, where symmetries can be exploited to reduce the search space~\cite{9617572}. Another application lies in data collections: for molecular databases, for example, it is sufficient to store the canonical form of a molecule rather than the molecule itself. Further applications can also be found in machine learning, software verification, and many other fields (as cited in~\cite{grohe2020graph}).

The main goal of this thesis is to re-implement one of the most efficient tools for graph isomorphism, \Naut{}, using a more modern programming language: C++. Recreating \Naut{} from scratch will allow us to evaluate optimisation differences, modernise its design, and extend its capabilities. For instance, the original \Naut{} does not support every possible graph family; in this updated version, we aim to include both some of the families already supported and additional ones, thereby contributing new features to the tool. The objective of this new codebase is certainly not to achieve better raw performance than the original \Naut{}, but rather to provide a modern perspective on the tool. Moreover, C++ remains a widely used programming language, and a C++ implementation of \Naut{} can make it more accessible to a broader audience while retaining better performance than more popular but interpreted languages such as Python.

This thesis is divided into four main chapters. The first introduces the fundamental concepts and preliminaries necessary to fully understand the key ideas discussed throughout this work. The second presents a comprehensive state of the art, summarising previous research on the graph isomorphism problem. The third is devoted to McKay's canonical labelling algorithm and the construction of the automorphism group. Finally, the fourth chapter presents our C++ implementation of \Naut{} and discusses the results obtained.